\documentclass[12pt, a4paper]{article}

%=============== PAQUETES ===============
\usepackage[utf8]{inputenc}
\usepackage[T1]{fontenc}
\usepackage[spanish]{babel}
\usepackage{geometry}
\usepackage{amsmath}
\usepackage{graphicx}
\usepackage{xcolor}
\usepackage{hyperref}
\usepackage{circuitikz}
\usepackage{comment}
\usetikzlibrary{arrows.meta, babel}

%=============== CONFIGURACIÓN ===============
\geometry{a4paper, margin=2.5cm}

\hypersetup{
    colorlinks=true,
    linkcolor=blue,
    urlcolor=cyan,
    pdftitle={Análisis Teórico: Regulador Zener},
}

%=============== DOCUMENTO ===============

\begin{document}

\title{Análisis Teórico de Circuitos con Diodos}
\author{Prof. César Rodríguez con asistencia de Gemini Code Assist}
\date{\today}
\maketitle

\section*{Análisis del Circuito Regulador de Tensión con Diodo Zener (Parte 1)}

Este documento presenta el análisis teórico del circuito regulador de tensión con diodo 
Zener propuesto en la 'Parte 1' del 'Procedimiento Experimental' de la guía de laboratorio 
`PRACTICA SEMANA 7`. El objetivo es determinar la tensión en la resistencia de carga 
$R_L$ ($V_{out}$) en función de la tensión de entrada $V_{in}$.

El circuito bajo análisis es el siguiente:

\begin{center}
\begin{circuitikz}[scale=1.2, transform shape] % Circuito mejorado para mayor claridad
    \draw (0,2) to[sV, l=$V_{in}$] (0,0) node[ground]{};
    \draw (0,2) to[R, l=$R_S$] (3,2);
    \draw (3,2) to[zD, *-*, l_=$D_Z$] (3,0) node[ground]{};
    \draw (3,2) to[R, l=$R_L$] (5,2) to[short] (5,0) node[ground]{};
    \draw (3,2) node[above] {$V_{out}$};
\end{circuitikz}
\end{center}

Donde:
\begin{itemize}
    \item $V_{in}$: Tensión de entrada.
    \item $R_S$: Resistencia en serie.
    \item $D_Z$: Diodo Zener con tensión de ruptura $V_Z$.
    \item $R_L$: Resistencia de carga.
    \item $V_{out}$: Tensión de salida a través de $R_L$ y $D_Z$.
\end{itemize}

\section*{Análisis de Operación}

El circuito opera en dos regiones principales, dependiendo del estado del diodo Zener.

\subsection*{Caso 1: Diodo Zener APAGADO (No en ruptura)}

Esta región ocurre cuando la tensión de entrada $V_{in}$ es baja y la tensión a través del diodo Zener es menor que su tensión de ruptura $V_Z$. En este caso, el diodo Zener se comporta como un circuito abierto (ignorando la pequeña corriente de fuga en inversa).

El circuito se simplifica a un divisor de tensión formado por $R_S$ y $R_L$. La tensión de salida $V_{out}$ se calcula como:

\begin{equation}
    V_{out} = V_{in} \cdot \frac{R_L}{R_S + R_L}
\end{equation}

El diodo Zener comenzará a conducir (entrará en ruptura) cuando la tensión $V_{out}$ alcance $V_Z$. Podemos encontrar la tensión de entrada mínima ($V_{in,th}$) para que el Zener se active:

\begin{equation}
    V_Z = V_{in,th} \cdot \frac{R_L}{R_S + R_L} \implies V_{in,th} = V_Z \cdot \frac{R_S + R_L}{R_L}
\end{equation}

Por lo tanto, esta región es válida para $V_{in} < V_{in,th}$.

\subsection*{Caso 2: Diodo Zener ENCENDIDO (En ruptura)}

Esta región se activa cuando la tensión de entrada $V_{in}$ es igual o superior a $V_{in,th}$. En este punto, el diodo Zener entra en su región de ruptura inversa y mantiene una tensión casi constante $V_Z$ a través de sus terminales.

En esta región, la tensión de salida $V_{out}$ se regula a la tensión Zener:

\begin{equation}
    V_{out} = V_Z
\end{equation}

Esta región es válida para $V_{in} \ge V_{in,th}$.

\section*{Conclusión: Tensión de Salida en Función de la Tensión de Entrada}
Combinando ambos casos, la tensión de salida $V_{out}$ en función de la tensión de entrada $V_{in}$ se puede expresar como una función por partes:

\begin{equation}
    V_{out}(V_{in}) = 
    \begin{cases} 
        V_{in} \cdot \frac{R_L}{R_S + R_L} & \text{si } V_{in} < V_Z \cdot \frac{R_S + R_L}{R_L} \\
        V_Z & \text{si } V_{in} \ge V_Z \cdot \frac{R_S + R_L}{R_L}
    \end{cases}
\end{equation}

Esta ecuación describe cómo el circuito regulador mantiene una tensión de salida constante $V_Z$ una vez que la tensión de entrada supera un umbral mínimo, y cómo se comporta como un simple divisor de tensión por debajo de ese umbral.

\newpage

\section{Análisis del Circuito Transitorio (Parte 2)}

Se analiza el circuito propuesto en la "Parte 2" de la guía, cuyo objetivo es estudiar el comportamiento transitorio de un diodo de señal.

El circuito es el siguiente:

\begin{center}
\begin{circuitikz}[scale=1.2, transform shape]
    % Fuente AC
    \draw (0,0) node[ground]{} 
        to[sV, l=$V_{in}$] (0,2) 
        
    % Resistencia (con formato de texto correcto para la unidad)
    % Usamos \text{k} para que no salga cursiva
        to[R, l=$R_1{=}1\text{k}\Omega$] (3,2) 
        
    % Diodo (etiqueta en modo texto normal es correcta)
        to[D, l=1N4148] (3,0) 
        node[ground]{};
        
    % Opcional: Si quieres visualizar la salida Vout sobre el diodo
    \draw (3,2) to[short, -o] (4,2) node[right] {$V_{out}$};
\end{circuitikz}
\end{center}

La tensión de salida $V_{out}$ es la tensión que cae sobre el diodo ($V_D$). Para determinar $V_{out}$ en función de $V_{in}$, utilizaremos el \textbf{modelo de caída de tensión constante} para el diodo. En este modelo, se asume que el diodo tiene una caída de tensión fija $V_{on}$ (típicamente 0.7V para diodos de silicio como el 1N4148) cuando está conduciendo.

\subsection*{Caso 1: Diodo APAGADO (Polarización Inversa)}

El diodo se comporta como un interruptor abierto cuando la tensión de entrada no es suficiente para polarizarlo en directa. Esto ocurre si $V_{in} \le V_{on}$.

En esta condición, no fluye corriente a través del circuito ($I_D = 0$). Al no haber corriente, no hay caída de tensión en la resistencia $R_1$. Aplicando la ley de tensiones de Kirchhoff (LVK) al lazo, la tensión de salida es igual a la de entrada.

\begin{equation}
    V_{out} = V_{in} \quad (\text{para } V_{in} \le V_{on})
\end{equation}

\subsection*{Caso 2: Diodo ENCENDIDO (Polarización Directa)}

El diodo se comporta como un interruptor cerrado con una fuente de tensión constante $V_{on}$ cuando la tensión de entrada supera este umbral. Esto ocurre si $V_{in} > V_{on}$.

En esta condición, el diodo "fija" la tensión en sus terminales a $V_{on}$, sin importar cuán grande sea $V_{in}$.

\begin{equation}
    V_{out} = V_{on} \approx 0.7\text{V} \quad (\text{para } V_{in} > V_{on})
\end{equation}

Este comportamiento es característico de un \textbf{circuito recortador (clipper)}, donde la tensión de salida sigue a la de entrada hasta un cierto nivel, y a partir de ahí se mantiene constante.

\subsection*{Conclusión: Función por Partes}
La relación entre $V_{out}$ y $V_{in}$ se describe como:
\begin{equation}
    V_{out}(V_{in}) = 
    \begin{cases} 
        V_{in} & \text{si } V_{in} \le V_{on} \\
        V_{on} & \text{si } V_{in} > V_{on}
    \end{cases}
\end{equation}

\newpage

\section{Análisis de Pequeña Señal (Parte 3)}

En esta sección se analiza el circuito de la "Parte 3", cuyo objetivo es verificar el concepto de \textbf{resistencia dinámica ($r_d$)} de un diodo. El circuito superpone una pequeña señal AC ($v_{ac}$) sobre una polarización DC ($V_{DC}$).

El circuito es el siguiente:

\begin{center}
\begin{circuitikz}[scale=1.2, transform shape]
    % Rama Izquierda: Fuentes en Serie
    % Dividí la altura total (3) en dos partes iguales (1.5 cada una)
    \draw (0,0) node[ground]{} 
        to[V, l=$V_{DC}{=}10\text{V}$] (0,1.5)  % Fuente DC (Batería)
        to[sV, l=$v_{ac}$] (0,3)                % Fuente AC (Sinusoidal)
        
    % Rama Superior: Resistencia
    % Usamos \text{k} para la unidad correcta
        to[R, l=$R_1{=}1\text{k}\Omega$] (3,3) 
        
    % Rama Derecha: Diodo
        to[D, l=1N4148] (3,0) 
        node[ground]{};

    % (Opcional) Terminal de salida para mayor claridad
    \draw (3,3) to[short, -o] (4,3) node[right] {$V_{out}$};
\end{circuitikz}
\end{center}

La tensión total de entrada es $V_{in}(t) = V_{DC} + v_{ac}(t)$. La tensión de salida $V_{out}(t)$ es la tensión en el diodo, $V_D(t)$. Para analizar este circuito, aplicamos el principio de \textbf{superposición}, separando el análisis en DC y AC.

\subsection{Análisis DC (Punto de Operación Q)}

Para el análisis DC, las fuentes AC se anulan (se convierten en cortocircuitos). El circuito se reduce a una fuente $V_{DC}$ en serie con la resistencia $R_1$ y el diodo. El objetivo es encontrar la corriente de polarización DC, $I_{DQ}$.

Usando el modelo de caída de tensión constante para el diodo ($V_D \approx V_{on} \approx 0.7$V), aplicamos la Ley de Voltajes de Kirchhoff (LVK):

\begin{equation}
    V_{DC} = I_{DQ} \cdot R_1 + V_{on}
\end{equation}

Despejando la corriente de polarización $I_{DQ}$:

\begin{equation}
    I_{DQ} = \frac{V_{DC} - V_{on}}{R_1}
\end{equation}

Esta corriente $I_{DQ}$ establece el \textbf{punto de operación (Q-point)} del diodo.

\subsection{Análisis AC (Pequeña Señal)}

Para el análisis AC, las fuentes DC se anulan (se convierten en cortocircuitos). El diodo se reemplaza por su modelo de pequeña señal: la \textbf{resistencia dinámica $r_d$}.

El valor de $r_d$ depende del punto de operación y se calcula como:
\begin{equation}
    r_d = \frac{V_T}{I_{DQ}}
\end{equation}
donde $V_T$ es la tensión térmica (aproximadamente 25mV a temperatura ambiente).

El circuito equivalente AC es un simple divisor de tensión formado por $R_1$ y $r_d$, alimentado por la fuente $v_{ac}$. La tensión de salida AC, $v_{out}$, es la tensión que cae en $r_d$:

\begin{equation}
    v_{out} = v_{ac} \cdot \frac{r_d}{R_1 + r_d}
\end{equation}

\subsection{Conclusión: Tensión de Salida Total}

La tensión de salida total $V_{out}(t)$ es la superposición de la componente DC y la componente AC:

\begin{equation}
    V_{out}(t) = V_{DQ} + v_{out}(t)
\end{equation}

Sustituyendo las expresiones encontradas:

\begin{equation}
    V_{out}(t) = V_{on} + v_{ac}(t) \cdot \frac{r_d}{R_1 + r_d}
\end{equation}

Donde $r_d$ se calcula usando la corriente $I_{DQ}$ del análisis DC. Esta ecuación muestra que la salida consiste en una tensión DC constante ($V_{on}$) sobre la cual se monta una versión \textbf{atenuada} de la señal AC de entrada. La atenuación depende de la relación entre $R_1$ y la resistencia dinámica $r_d$, que a su vez depende de la polarización DC ($V_{DC}$).

\end{document}
